\family{Combretaceae}{Terminalia family}
\label{fam:combretaceae}

Trees, shrubs or climbers. Leaves simple and entire, alternate or
opposite, in \taxon{Terminalia} crowded at the ends of short shoots on
tiered branches. Flowers small, in spikes or heads. Fruit dry and
indehiscent, usually with two to five wings. Two species of
\taxon{Terminalia} in the Arami Uplands, both important timber trees.

\section*{Key to the species}
\begin{florakey}
	\lead{1a}{Bark dark and deeply fissured; buttresses small; fruit
		longer than broad}{\taxon{Terminalia ivorensis}}
	\lead{1b}{Bark pale grey and smooth; tall plank buttresses; fruit much
		broader than long}{\taxon{Terminalia superba}}
\end{florakey}

\species{Terminalia ivorensis}{A.Chev.}
\label{sp:COM01}

\field{Local and trade names} Yoruba: Idigbo. Trade: Idigbo.

\field{Status} Native.

\field{Description} Deciduous tree to 45~m with a straight,
cylindrical bole and small buttresses. Bark dark grey to black, deeply
fissured. Branches in tiers. Leaves crowded at the ends of the shoots,
elliptic to obovate, 5--10~cm long. Flowers small, cream, in spikes.
Fruit with two wings, oblong, longer than broad.

\field{Phenology} Flowering February to April; fruiting April to June.

\field{Habitat} Lowland semi-deciduous forest.

\field{Uses} A light, easily worked timber used for joinery, window
frames and furniture; also planted in forest plantations
\citep{louppe2008}.

\field{Specimens examined} Igbo-Oba Forest Reserve (C3), 470~m, st., 3~Sep 1935, \specimen{Harding}{1378}{OKF};
Ireke Ferry (D4), 360~m, fr., 15~Jun 1936, \specimen{Harding}{1483}{OKF};
Igbo-Oba Forest Reserve (C3), 470~m, fr., 16~May 1937, \specimen{Harding}{1696}{OKF};
Ireke Ferry (D4), 360~m, fr., 11~Apr 1938, \specimen{Harding}{1892}{OKF};
Igbo-Oba Forest Reserve (C3), 470~m, st., 8~Dec 2020, \specimen{Alade}{399}{AUH};
Ireke Ferry (D4), 360~m, st., 11~Jan 2021, \specimen{Alade}{435}{AUH};
Okeri (C4), 420~m, st., 22~Aug 2023, \specimen{Alade}{806}{AUH}.

\field{Distribution} Grid squares C3, C4, D4.

\begin{center}
\distmap{C/3, C/4, D/4}
\end{center}

\species{Terminalia superba}{Engl. \& Diels}
\label{sp:COM02}

\field{Local and trade names} Yoruba: Afara. Trade: Afara.

\field{Status} Native.

\field{Description} Deciduous tree to 50~m with tall plank buttresses.
Bark smooth, pale grey. Branches in tiers. Leaves crowded at the ends
of the shoots, obovate, 6--12~cm long. Flowers small, greenish-white,
in spikes. Fruit with two wings, much broader than long.

\field{Phenology} Flowering February to April; fruiting May to August.

\field{Habitat} Lowland semi-deciduous forest and forest--savanna
mosaic, often in secondary forest.

\field{Uses} A pale timber used for plywood, furniture and interior
work, and widely planted in forest plantations
\citep{keay1989,louppe2008}.

\field{Specimens examined} Agbede (C2), 540~m, fr., 20~Jun 1934, \specimen{Harding}{1139}{OKF};
Ireke Ferry (D4), 360~m, fr., 16~May 1937, \specimen{Harding}{1703}{OKF};
Igbo-Oba Forest Reserve (C3), 470~m, fr., 23~Aug 1938, \specimen{Harding}{1992}{OKF};
Ewe Stream gallery forest (B3), 580~m, fr., 17~Jun 1962, \specimen{Olaniyan}{413}{SVN};
Aramiwo Foot (B2), 640~m, fr., 22~Jun 1966, \specimen{Olaniyan}{805}{SVN};
Ewe Stream gallery forest (B3), 580~m, fr., 26~Aug 2021, \specimen{Alade}{544}{AUH};
Oke-Asa farmland (B4), 690~m, fr., 5~Jun 2023, \specimen{Alade}{777}{AUH}.

\field{Distribution} Grid squares B2, B3, B4, C2, C3, D4.

\begin{center}
\distmap{B/2, B/3, B/4, C/2, C/3, D/4}
\end{center}
